\documentclass[11pt,a4paper]{appletmanual}
\usepackage[spanish,es-noquoting,es-nodecimaldot]{babel}
\manuallang{es}

\appletname{Restricciones No Lineales}
\appletsub{Siete esquemas de subsidio doblan la recta de presupuesto en sitios
distintos. La pregunta es siempre la misma: dónde cae la mejor canasta
asequible, y si cae justo en el pliegue.}
\appletdir{nonlinear-budget}

\begin{document}
\manualtitlepage{Manual del applet · Microeconomía I · Universidad de los Andes\\
Traducido del applet de GeoGebra «Max. Utilidad — Rest. Presupuestaria No Lineal».\\
Código MIT. Se abre en el navegador; no hay nada que instalar.}

\manualtoc
\thispagestyle{empty}
\clearpage
\setcounter{page}{1}

% ==========================================================================
\section{Qué es esto}

La restricción presupuestaria de un manual es una recta. Casi ningún programa
público deja una recta. Un cupo de leche gratis, un descuento sobre las
primeras unidades, una plaza en un colegio público que se pierde si se compra
una privada: cada uno dobla la frontera en un sitio distinto, y el óptimo puede
caer exactamente en el pliegue.

Este applet trae siete esquemas, la posibilidad de comparar dos a la vez, y un
óptimo calculado de manera que sí encuentra los pliegues.

\figher[1.0]{kink}{Subsidio en especie no excluyente. Las primeras $x_s$
unidades son gratis, y a partir de ahí se paga el precio de mercado: la
frontera tiene un tramo horizontal y luego un quiebre.}

\begin{amkey}
La comparación que importa es entre un subsidio en especie y la transferencia
en efectivo que cuesta lo mismo a quien la financia. Si el cupo no restringe,
las dos empatan. Si restringe, el efectivo gana, y el applet mide por cuánto.
\end{amkey}

\subsection{Para quién}

Para el tema de restricciones presupuestarias, y para cualquier clase de
economía pública donde se discuta transferencias en especie frente a efectivo.
Supone curvas de indiferencia y la idea de óptimo; no supone cálculo.

%% Secciones comunes a todos los manuales, edición española.
%% Se incluyen con \input{common-es}. Lo único que cambia entre applets es
%% \appletfolder, que es también el nombre de la carpeta y de la ruta en el sitio.

\section{Cómo abrirlo}

No hay que instalar nada. La aplicación entera es un solo archivo
\code{index.html}: no descarga librerías, no llama a ningún servidor y no
guarda nada fuera del navegador.

\subsection{Las tres maneras}

\begin{description}
\item[Doble clic sobre el archivo.] Es lo más rápido. Busque
  \code{\appletfolder/index.html} y ábralo. El navegador lo muestra desde el disco,
  con la barra de direcciones empezando por \code{file://}. Funciona sin
  conexión desde el primer momento.
\item[Desde la web.] Si alguien ya lo publicó, la dirección termina en
  \code{/\appletfolder/}. Una vez cargada la página, puede desconectarse: no
  necesita la red para nada más.
\item[Servido desde una carpeta.] Para una clase, poner la carpeta en
  cualquier servidor estático basta. Con Python instalado:
  \begin{quote}\code{python3 -m http.server 8000}\end{quote}
  y luego \code{http://localhost:8000/\appletfolder/} en el navegador.
\end{description}

\begin{amnote}
Guardar la página con \key{Ctrl}\,+\,\key{S} (o \key{Cmd}\,+\,\key{S} en
Mac) deja una copia propia que sigue funcionando sin conexión y que se puede
llevar en una memoria USB o repartir por correo.
\end{amnote}

\subsection{Idioma y tema}

Arriba a la derecha hay dos pares de botones.

\begin{ctrltable}{Control & Qué hace}
ES / EN & Cambia el idioma de toda la interfaz al vuelo. Nada se recalcula:
          puede cambiarlo en mitad de una explicación sin perder el estado. \\
Oscuro / Claro & Cambia el tema. El oscuro es el diseño original; el claro
          existe para proyectores y para imprimir. \\
\end{ctrltable}

Las dos elecciones se recuerdan en el navegador, así que la próxima vez se abre
como la dejó.

\subsection{Qué navegador}

Cualquiera que sea razonablemente reciente: Chrome, Edge, Firefox o Safari, en
ordenador, tableta o teléfono. En pantallas estrechas los paneles se apilan uno
sobre otro en lugar de ponerse al lado. No hace falta cuenta, ni permisos, ni
extensiones.


% ==========================================================================
\section{Los siete esquemas}

Con precio $p_x$, renta $m$, cupo $x_s$ y tasa de subsidio $s$:

\begin{ctrltable}{Núm. y nombre & Frontera}
0 · Restricción simple & La recta de siempre. \\
1 · Subsidio en especie, no excluyente & Plana hasta $x_s$, y luego al precio
  de mercado: un quiebre. \\
2 · Subsidio parcial, no excluyente & Precio $(1-s)p_x$ hasta $x_s$, luego
  $p_x$: un quiebre más suave. \\
3 · Subsidio excluyente (colegios públicos) & O se toman $x_s$ gratis, o se
  compra al precio $p_x$: una caída vertical. \\
4 · Subsidio parcial excluyente & Descuento sobre las primeras $x_s$ unidades,
  que se pierde entero si se pasa del cupo. \\
5 · Equivalente en efectivo (total) & Renta $m + p_x x_s$, recta. \\
6 · Equivalente en efectivo (parcial) & Renta $m + s\,p_x x_s$, recta. \\
\end{ctrltable}

El 1 se empareja con el 5, y el 2 con el 6: el mismo coste para quien financia
el programa, sin condiciones. Poner uno frente al otro es para lo que sirve el
control \ui{Comparar con}, y es el asunto del applet.

\begin{amnote}
Con $s < 0$ el subsidio se convierte en un impuesto de tipo $\tau = -s$ sobre
las primeras $x_s$ unidades. El deslizador llega hasta $-1$, así que los siete
esquemas sirven también para hablar de impuestos selectivos.
\end{amnote}

\subsection{Quiebre contra salto}

Los esquemas 1 y 2 producen un \textbf{quiebre}: la frontera cambia de
pendiente pero no se corta. Los esquemas 3 y 4 producen un \textbf{salto}: pasar
del cupo hace perder el subsidio entero, así que la frontera cae en vertical y
hay canastas que ningún gasto alcanza.

\fig[1.0]{drop}{Subsidio excluyente. La frontera cae en vertical en $x_s$: una
unidad más de $x$ cuesta perder el cupo entero.}

% ==========================================================================
\section{La pantalla}

\fig[1.0]{interface}{La pantalla completa.}

\begin{description}
\item[Esquema.] Los siete, numerados de 0 a 6.
\item[Comparar con.] Dibuja un segundo esquema detrás del primero, y activa la
  lectura \ui{Ventaja del efectivo}.
\item[Preferencias.] Siete familias, la última escribible.
\item[Precios y renta.] $p_x$, $p_y$, $m$.
\item[El programa.] El cupo $x_s$ y la tasa $s$.
\item[Vista.] Degradado de calor, panel de sustitución, ejes.
\item[Capas.] Qué se dibuja.
\end{description}

\subsection{Los dos paneles}

\begin{description}
\item[Plano $(x, y)$.] El mapa de utilidad bajo una rampa topográfica, lo
  inasequible velado, la frontera del esquema elegido, la del comparado detrás,
  la curva de indiferencia por el óptimo y los vértices.
\item[Utilidad sobre la restricción.] $u(x, f(x))$ conforme $x$ recorre la
  frontera. Su punto más alto \emph{es} el óptimo, así que «encontrar la mejor
  canasta asequible» se convierte en «encontrar la cima de esta curva». En los
  esquemas 3 y 4 la curva se parte visiblemente en dos: ese hueco es el
  subsidio que se pierde.
\end{description}

\fig[0.62]{drop-path}{La utilidad sobre la restricción en un esquema
excluyente. El hueco entre los dos arcos es exactamente el subsidio perdido.}

\subsection{El tema oscuro}

\fig[0.95]{dark}{El plano y el panel de sustitución en el tema oscuro.}

% ==========================================================================
\section{Cómo se encuentra el óptimo}

Una condición de primer orden es el instrumento equivocado aquí. En estas
fronteras la mejor canasta está muy a menudo en un vértice, donde la pendiente
no está definida y no hay tangencia que resolver; y en los esquemas excluyentes
está en el extremo derecho de la primera pieza, con un punto estrictamente peor
inmediatamente a su derecha.

\begin{enumerate}
\item Cada pieza de la frontera se barre densamente.
\item La mejor muestra interior de cada pieza se refina por sección áurea.
\item Se prueba \emph{cada vértice explícitamente}, por los dos lados de cada
  unión.
\item Gana el mejor de todos. De dónde salió es lo que lo clasifica en
  tangencia, vértice, esquina o borde del salto.
\end{enumerate}

\begin{ctrltable}{Tipo de óptimo & Qué significa}
Tangencia & Interior a un tramo: la RMS iguala al precio relativo vigente ahí.
  El pliegue no muerde. \\
Vértice & Justo en el codo. A la izquierda la restricción es demasiado plana y
  a la derecha demasiado inclinada: lo mejor es quedarse en el pliegue. El
  consumidor agota el cupo exacto. \\
Esquina & Sobre un eje: conviene gastarlo todo en un bien. \\
Borde del salto & En un esquema excluyente, pegado al borde. Una unidad más
  dejaría al consumidor estrictamente peor. \\
\end{ctrltable}

% ==========================================================================
\section{Las preferencias}

\begin{ctrltable}{Familia & Forma}
Cobb--Douglas & $x^{a} y^{1-a}$ \\
Sustitutos perfectos & $a x + (1-a) y$ \\
Complementarios perfectos & $\min(a x, (1-a) y)$ \\
Cuasilineal & $8a\ln x + y$ \\
CES generalizada & $\bigl[\tfrac{1}{2a} + \tfrac{1}{2(1-a)}\bigr]^{-1}
   (a x^{\rho} + (1-a) y^{\rho})^{1/\rho}$ \\
X inferior & $\ln(x - \underline{x}) - \tfrac{1}{1-\beta}\ln(\bar y - y)$ \\
Personalizada & Lo que escriba \\
\end{ctrltable}

Un solo peso $a$ recorre todas: el exponente en la Cobb--Douglas, el valor
relativo en el caso lineal y en el de Leontief, la participación en la CES.

\begin{amnote}
Dos notas sobre la CES generalizada. Su factor de cabecera se escribe
internamente como $2a(1-a)$, que es idénticamente igual: el literal
$1/(2a) + 1/(2(1-a))$ divide por cero en los dos extremos del deslizador,
mientras que $2a(1-a)$ simplemente tiende a cero ahí. Siendo una constante
positiva, reescala la utilidad sin tocar las preferencias ni el óptimo. Y en
$\rho = 0$ la CES no está definida aunque su límite sea la Cobb--Douglas, así
que un deslizador parado en cero emite ese límite en vez de dividir por él.
\end{amnote}

\fig[1.0]{complements}{Complementarios perfectos sobre un subsidio en especie.
Con preferencias de Leontief el óptimo está siempre en el codo de la curva de
indiferencia, y aquí además compite con el codo de la frontera.}

% ==========================================================================
\section{Los controles}

\begin{ctrltable}{Deslizador & Qué es}
$p_x$, $p_y$ & Precios, acotados por debajo en $0{,}2$. \\
$m$ & Renta. \\
$x_s$ & El cupo del programa. Se recorta a lo que la renta puede comprar de
  verdad, para que el subsidio nunca supere al presupuesto. \\
$s$ & La tasa de subsidio, de $-1$ a $1$. Negativa es un impuesto. \\
$a$, $\rho$, $\beta$, $\underline{x}$, $\bar y$ & Parámetros de la familia de
  preferencias elegida; solo aparecen los que esa familia usa. \\
\end{ctrltable}

\subsection{Vista}

\begin{description}
\item[Degradado de calor.] Apagado, el fondo queda liso. Las curvas de
  indiferencia conservan su color de altura en todo momento: con el mapa puesto
  concuerda con el suelo, y con el mapa quitado es la única lectura de altura
  que queda.
\item[Panel de sustitución.] La lectura de utilidad sobre la restricción es una
  vista de apoyo, así que ocupa menos ancho y se puede quitar del todo para
  dejarle el marco entero al plano.
\item[Ejes.] Fijos por defecto, puestos con las casillas $x \le$ e $y \le$. Un
  eje que se reajusta en cada movimiento renumera mientras se arrastra, y
  entonces dos ajustes no se pueden comparar a ojo —que es justo el objetivo de
  poner un esquema y su equivalente en efectivo en el mismo dibujo. El ajuste
  automático sigue disponible, y si una frontera se sale del marco fijo aparece
  un aviso en vez de recortarla en silencio.
\end{description}

\subsection{Capas}

\begin{ctrltable}{Capa & Qué dibuja}
Curvas de fondo & Curvas de nivel de la utilidad. \\
Conjunto factible & Vela lo inasequible. \\
Curva por el óptimo & La curva de indiferencia que pasa por él. \\
Óptimo & El punto elegido. \\
Vértices & Los codos y bordes de la frontera. \\
Conjunto comparado & La frontera del esquema con el que se compara. \\
Retícula & La cuadrícula de fondo. \\
\end{ctrltable}

% ==========================================================================
\section{Actividades para clase}

\begin{activity}{La recta de siempre}
\amset Esquema \ui{0 · Restricción simple}, valores de partida.
\amexp $x^\ast = 0{,}69$, $y^\ast = 2{,}18$, utilidad $1{,}539$, coste $0{,}00$.
El tipo de óptimo es \ui{Tangencia}: el caso del manual.
\par\smallskip
Es la línea de base contra la que se leen los seis esquemas siguientes.
Anótela en la pizarra antes de seguir.
\end{activity}

\begin{activity}{Cuando el cupo no muerde}
\amset Esquema \ui{1 · Subsidio en especie, no excluyente}, valores de partida.
Ponga \ui{Comparar con} en \ui{5 · Equivalente en efectivo (total)}.
\amexp $x^\ast = 0{,}96$, $y^\ast = 3{,}05$, utilidad $2{,}158$, coste
$3{,}50$. La \ui{Ventaja del efectivo} da $0{,}000$ y el veredicto dice
«Efectivo y especie empatan».
\par\smallskip
Mire por qué: el óptimo cae en el tramo donde las dos fronteras coinciden. El
cupo $x_s = 0{,}92$ es más pequeño de lo que este consumidor iba a comprar de
todas formas, así que el programa le da dinero y nada más. Un subsidio en
especie que no restringe \emph{es} una transferencia en efectivo.
\end{activity}

\begin{activity}{Cuando el cupo sí muerde}
\amset Igual que la anterior, pero suba $x_s$ a $2{,}20$.
\amexp $x^\ast = 2{,}20$ exactamente —el cupo entero, ni una unidad más—,
$y^\ast = 3{,}11$, utilidad $2{,}801$, coste $8{,}36$. El tipo de óptimo pasa a
\ui{Vértice} y la \ui{Ventaja del efectivo} sube a $+0{,}217$.
\par\smallskip
Con el mismo coste de $8{,}36$, el efectivo alcanza una curva de indiferencia
más alta. Pregunte por qué un gobierno elegiría entonces la especie. Las
respuestas —paternalismo, focalización, externalidades del bien subsidiado— son
justo las que no aparecen en este modelo, y por eso conviene sacarlas aquí.
\end{activity}

\begin{activity}{Un salto en lugar de un quiebre}
\amset Esquema \ui{3 · Subsidio excluyente (colegios públicos)}, valores de
partida.
\amexp La frontera cae en vertical en $x_s$: la lectura \ui{Salto en $x_s$} da
$-1{,}25$. El tipo de óptimo es \ui{Borde del salto}.
\par\smallskip
Encienda el \ui{Panel de sustitución}. La curva de utilidad sobre la
restricción se parte en dos arcos que no se tocan, y el hueco entre ellos es
exactamente el subsidio que se pierde al cruzar. Pregunte qué le pasa a una
familia que necesita apenas una unidad más que el cupo.
\end{activity}

\begin{activity}{Impuesto en vez de subsidio}
\amset Esquema \ui{2 · Subsidio parcial, no excluyente}. Baje $s$ hasta $-0{,}6$.
\amexp El pliegue cambia de sentido: ahora las primeras $x_s$ unidades cuestan
más caras y la frontera es \emph{cóncava} en lugar de convexa. El óptimo se
aleja de $x$.
\par\smallskip
Es un impuesto selectivo sobre las primeras unidades, y el mismo dibujo sirve
para explicarlo. Pregunte qué esquema fiscal real se parece a esto.
\end{activity}

\begin{activity}{Preferencias que siempre acaban en el codo}
\amset Preferencias \ui{Complementarios perfectos}, esquema
\ui{1 · Subsidio en especie}.
\amexp Con Leontief el óptimo está siempre en el codo de la curva de
indiferencia. Mueva $x_s$: mientras el codo del consumidor esté a la derecha
del cupo, el subsidio se aprovecha entero y luego se sigue comprando.
\par\smallskip
Compárelo con \ui{Sustitutos perfectos}, donde el óptimo salta de un eje al
otro en cuanto el precio relativo cruza la relación marginal de sustitución.
Los dos casos extremos hacen de tope a la intuición.
\end{activity}

\begin{activity}{Los ejes fijos no son un capricho}
\amset Cualquier esquema con \ui{Comparar con} activo. Marque y desmarque
\ui{Ajustar los ejes solos} mientras mueve $m$.
\amexp Con los ejes automáticos, los números del eje cambian a cada paso y las
dos fronteras parecen quedarse quietas. Con los ejes fijos se ve cuál se mueve
y cuánto.
\par\smallskip
Es el motivo de que los ejes vengan fijos por defecto: la comparación es el
objeto del applet, y una comparación cuya escala cambia sola no es una
comparación.
\end{activity}

\begin{activity}{Construir el caso en que la especie gana}
\amset Libre.
\amexp No se puede, y ese es el ejercicio. Con las preferencias que sean, el
efectivo del mismo coste alcanza siempre una curva de indiferencia igual o más
alta, porque el conjunto factible del efectivo contiene al de la especie.
\par\smallskip
Lo más que se consigue es empatar, y se empata exactamente cuando el cupo no
restringe. Que la clase lo intente durante cinco minutos vale más que
enunciarlo.
\end{activity}

% ==========================================================================
\section{Sintaxis de expresiones}

\begin{ctrltable}{Elemento & Qué se admite}
Operadores & \code{+ - * / \^{} ( )} \\
Funciones & \code{sqrt ln log log10 exp abs sin cos tan min max} \\
Constantes & \code{e}, \code{pi} \\
Multiplicación implícita & \code{2xy} es \code{2*x*y} \\
Asociatividad & \code{\^{}} liga por la derecha y más fuerte que el menos
  unario: \code{-x\^{}2} es $-(x^2)$ \\
\end{ctrltable}

% ==========================================================================
\section{Diferencias con el original de GeoGebra}

\begin{description}
\item[Las discontinuidades están modeladas, no aproximadas.] El original
  escribía los esquemas 3 y 4 como una sola función \code{If[...]}, que GeoGebra
  dibuja con un segmento que une los dos lados del salto. Aquí cada esquema es
  una lista de piezas separadas, así que la caída sigue siendo una caída y el
  perfil de utilidad se parte donde debe.
\item[Rango de color.] \code{uMIN := U(0.001, 0.001)} no está definida para la
  utilidad de Stone-Geary, cuyo dominio exige $x > \underline{x}$; la rampa se
  quedaba en un solo color para todo el ejemplo de «x inferior». La reescritura
  toma el percentil 2 del campo muestreado.
\item[\code{uMAX}] usaba el mismo valor derivado de $x$ para las dos
  coordenadas, y superaba $\bar y$ lo bastante a menudo como para devolver
  \code{NaN}. Misma corrección.
\item[Precios.] $p_X$ y $p_Y$ empezaban en 0, con lo que $I/p_Y$ y toda la
  construcción dividían por cero. Los dos están acotados por debajo en $0{,}2$.
\item[$x_s$ y $\underline{x}$] se recortan a lo que la renta puede comprar, así
  que el cupo nunca supera al presupuesto y la utilidad de Stone-Geary sigue
  definida.
\item[Objetos muertos.] \code{pXlow} y \code{pXhigh} se calculaban con una
  fórmula que devuelve números negativos grandes para todos los ajustes que
  permiten los deslizadores, y no se mostraban en ninguna parte. No se
  trasladaron.
\item[La comparación con efectivo se calcula], no se deja para que la estime el
  lector: el coste de cada esquema y la diferencia de utilidad contra la
  transferencia en efectivo del mismo coste son lecturas.
\end{description}

% ==========================================================================
\section{Problemas frecuentes}

\begin{ctrltable}{Síntoma & Qué pasa}
Aviso de que la frontera se sale & Los ejes están fijos y el dibujo no cabe.
  Suba $x \le$ o $y \le$, o marque el ajuste automático. \\
El plano sale liso & El \ui{Degradado de calor} está apagado. Las curvas
  conservan su color de altura igualmente. \\
La ventaja del efectivo dice «---» & No hay esquema en \ui{Comparar con}. \\
El óptimo no se mueve al cambiar $s$ & El cupo no está restringiendo: el óptimo
  cae fuera del tramo subsidiado. Suba $x_s$. \\
La utilidad sale \code{NaN} & En la familia «X inferior» el dominio exige
  $x > \underline{x}$ e $y < \bar y$. Ajuste esos dos deslizadores. \\
\end{ctrltable}

%% Secciones comunes a todos los manuales, edición española.
%% Se incluyen con \input{common-es}. Lo único que cambia entre applets es
%% \appletfolder, que es también el nombre de la carpeta y de la ruta en el sitio.

\section{Publicarlo para una clase}

Sirve cualquier alojamiento estático: GitHub Pages, Netlify, un directorio web
de la universidad, o incluso una carpeta compartida servida por HTTP. Lo único
que hay que subir es el archivo \code{index.html} dentro de una carpeta con el
nombre que quiera que aparezca en la dirección.

El repositorio trae un flujo de trabajo de GitHub Actions
(\code{.github/workflows/pages.yml}) que publica todos los applets a la vez en
cada empuje a la rama principal. Para activarlo, en el repositorio:
\emph{Settings\,$\rightarrow$\,Pages\,$\rightarrow$\,Source: GitHub Actions}.
Es lo único que no puede hacer el propio flujo, porque crear un sitio de Pages
requiere permisos de administración que un \code{GITHUB\_TOKEN} nunca recibe.

\begin{amnote}
Para repartir el applet a estudiantes sin conexión fiable, mándeles el archivo
\code{index.html} directamente. Pesa menos que una fotografía y se abre con
doble clic.
\end{amnote}

\section{Licencia y procedencia}

El código de la aplicación es MIT. Puede usarlo, modificarlo y repartirlo, en
clase o fuera de ella, citando la procedencia.

Este applet es una reescritura autónoma de un applet original de GeoGebra. La
reescritura no es una traducción literal: donde el original tenía un error de
construcción, una división por cero alcanzable con los deslizadores, o un objeto
muerto heredado del applet del que se copió, esta versión lo corrige y lo
documenta. La sección \emph{Diferencias con el original de GeoGebra} enumera
esos cambios uno por uno.

Todo está escrito a mano y sin dependencias: el analizador de expresiones, la
derivación simbólica, el trazado de curvas de nivel y el dibujo. En particular
no se usa \code{eval}, de modo que lo que escribe un estudiante en una casilla
de texto no puede convertirse en código, y una política de seguridad de
contenidos estricta no rompe la página.


\end{document}
